\section{Expansion gates}\label{sec:05-gates}

We show that a sufficiently long carried arc expands in one direction,
and that enough expansion determines a unique continuation in that
direction. These path lemmas will be applied to the ordinary rows and,
later, to relative rows.  The relative statement is conditional on the full
package below; passing to relative rows does not automatically preserve that
package.  We retain the ambient torsion-free hyperbolic group, its height
homomorphism, and the cyclic noncommensuration property established above.
In particular, a nonzero-height element cannot commensurate a nontrivial
height-zero cyclic subgroup.

\subsection{Carried arcs and bounded comparisons}

For the ordinary application, sample successive coarse Magnus edges
of \(T_0\). The short window is the sampled edge itself, its group
\(H_i\) is the corresponding free Magnus group, and \(R_i\) is its
free Cayley tree. There are no peripheral groups in this application.
The row geometry and path-inertia results supply quasiconvex window
groups, finite orbit lists for each bounded shape, and the cyclic
intersection bound for competing continuations. We formulate the
argument more generally because Section~\ref{sec:07-transfer} will
use relative free-product trees in place of these Cayley trees.

\begin{definition}[The expansion package]\label{gate:package}
Consider oriented monotone paths in a tree of groups with height.  Sample at
successive occurrences of one edge type, and number these occurrences by
integers.  At position $i$ fix a short window $W_i$, of a fixed radius $s$;
the window may consist of the sampled edge itself.  Write
$H_i=\Stab(W_i)$, with pointwise stabilizers throughout.  In the relative
case fix an invariant system of height-zero peripheral groups.  A group is
\emph{nonparabolic} if it is nontrivial and is not contained in any one of
these actual peripheral groups.  The required package consists of the
following properties.
\begin{enumerate}[label=(E\arabic*)]
\item The total tree of ordinary or relative Cayley spaces is hyperbolic,
and all adjacent edge maps are uniformly quasi-isometric embeddings.
\item Nonparabolically stabilized windows of each fixed shape have finitely
many orbits.  The same holds for their nonparabolically stabilized
extensions, including configurations with a specified common subwindow.
Their stabilizers, with intrinsic relative metrics, embed quasi-isometrically
in every row of the window; constants are uniform for each bounded shape.
\item A stabilizer of a window containing a short window is a free product
of representatives of its nontrivial peripheral intersections and a free
group, with a finite relative decomposition.  The peripheral systems
restrict exactly under further window restrictions.  Every window group
under consideration has a loxodromic element in the relevant short-window
free-product trees.
\item Two distinct continuations through the same short window, with their
terminal positions at the same height, have stabilizers intersecting in an
at most cyclic group.
\end{enumerate}
For $H_i$ use the Bass tree $R_i$ of the stated free-product decomposition,
with trivial edge stabilizers and finite quotient.  Its elliptic subgroups
are exactly the prescribed restricted peripheral subgroups.  In the ordinary
application $R_i$ is the horizontal free Cayley tree, and the deep-window
groups under consideration are noncyclic.  Relative trees may have vertices
of infinite valence.
\end{definition}

A \emph{trajectory} is one of these monotone paths through the row tree;
it may be finite, one-sided infinite, or bilateral. The arcs carried
along it lie in the corresponding trees $R_i$.

The choices in this definition, including the comparisons constructed next,
are equivariant under translation of a path together with its marked
positions.  For a group acting on $R_i$, its \emph{core} means its minimal
invariant subtree; its quotient will be called the \emph{core graph}.
These are different objects: the subtree is generally infinite, while
the core graphs below have finitely many edges. An edge of \(R_i\)
will also be called a \emph{regular edge}, to emphasize its trivial
stabilizer. An arc is \emph{carried} by a window if it lies on a
bi-infinite line in that window group's core.  For an infinite path, carrying
means carrying by every relevant finite subwindow.

\begin{lemma}[Bounded comparisons]\label{gate:comparisons}
Assume the expansion package.  For a fixed bounded pair of positions $i,j$,
let $K$ stabilize the segment spanned by $W_i\cup W_j$.  There are
$K$-equivariant coarse quasi-isometries between the $K$-cores in $R_i,R_j$.
They have coarse inverses, preserve the line endspaces of every further
window, and are coarsely compatible on every fixed finite diagram of
positions.  The constants depend on the bounded diagram, not on a further
window carrying the arc.
\end{lemma}
\begin{proof}
The $K$-action on either core is cocompact.  To see this, take the finite
relative generating set supplied by (E3), paths from a base vertex to its
translates by those generators, and paths to fixed vertices of the finitely
many peripheral factors.  The translates of the resulting finite subtree
contain the minimal subtree.  Conversely, a segment in the core reads as a
bounded-cost word in relative generators: there are finitely many edge
representatives, and changes between their lifts are accounted for by the
vertex stabilizers.  Exact restriction of the elliptic systems identifies
these vertex stabilizers with the allowed peripheral moves.  Thus the orbit
metrics on both cores are quasi-isometric to the same intrinsic relative
Cayley metric of $K$.

For each representative vertex choose an image fixed by its stabilizer in
the other tree; such a vertex exists by exact ellipticity.  Extend these
choices equivariantly and then over edges.  Finite quotient gives bounded
choices, coarse inverses, and the asserted quasi-isometries.  For a fixed
finite diagram, first compare the maps and their composites on the core
of the common window spanned by that diagram.  They are equivariant for
its group and agree within bounded distance on a finite set of orbit
representatives.  Property (E2) makes this bound uniform for the bounded
shape of the diagram.  A further window has its core inside this common
core, so restriction inherits the same bound; no coboundedness constant
depending on that further window is introduced.  This proves compatibility
and preservation of the boundary points of further window cores.

We use the following precise convention for arcs.  Join the images of the
endpoints by a geodesic and delete a fixed bounded segment at each end.
The remaining arc lies in every transferred line containing the original
arc.  The trimming constant follows from stability of quasi-geodesics in a
tree.  Under immediate inverse comparison, the transferred arc is contained
in a bounded enlargement of the original arc; in particular its length is
at most the original length plus a uniform additive constant.
\end{proof}

\subsection{The corridor estimate and phase uniqueness}

\begin{lemma}[Flaring]\label{gate:flaring}
Under the expansion package there are integers $b,R>0$ with the following
property.  An arc of length at least $R$ carried across the positions
$[-b,b]$ has trimmed image of length at least $100$ times its original
length in one of the comparisons $0\to-b$ and $0\to b$.
\end{lemma}
\begin{proof}
We establish the corridor estimate with its dependence on the window
length made explicit.  Fix $m$.  Work with the stabilizer of the comparison
window through $[-m,m]$, including its short tails.  Its core is cobounded,
so the endpoints of a carried arc at position $0$ can be approximated by
two orbit points.  The error may depend on $m$ but not on the arc.  Translate
one fixed section through this window by the corresponding two group
elements.  These sections form the two sides of a quadrilateral and cross
each sampled fiber exactly once.  Their total length, and all endpoint
approximation errors, are bounded in terms of $m$, uniformly over window
types.

Let $l_j$ be the distance between the untrimmed comparison images at
position $j$, with $l_0$ the original length.  Use the fixed comparisons
through the short windows, rather than comparisons chosen afresh for a
larger carrying window.  By Lemma~\ref{gate:comparisons}, the horizontal
width between the sections is comparable to $l_j$ with uniform
multiplicative constants and an additive error depending on $m$.

For $d\le m$, fill the loop formed by the two sections and horizontal
paths in the rows $-d,d$ by a linear-area Rips disc in the hyperbolic total
space \cite[Proposition~III.H.2.7]{BridsonHaefliger1999}.
Here the Rips parameter and the linear isoperimetric constant are
fixed.  We explain why its area bounds the sum of the intervening fiber
widths.  Extend projection to the underlying tree across the edges of the
triangulation by bounded paths and across triangles by tripod fillings.
A transverse cut through an intervening tree edge has two boundary
crossings, one on each section.  A component of its preimage joins those
crossings: preimage tracks pair boundary crossings, and these are the only
two.  A small perturbation arranges transversality without altering the
boundary projection.

Label consecutive track crossings by points of the relevant actual fiber
within bounded total-space distance of the triangle vertices.  Such points
exist because a bounded path whose projected endpoints lie on opposite
sides crosses the fiber.  Two resulting nearby fiber points have uniformly
bounded intrinsic distance.  Indeed a bounded total-space path between
them has a bounded projected edge word; successively remove its backtracks,
using the uniform quasi-isometric edge embeddings.  This bounds the
intrinsic fiber displacement by a constant depending only on the fixed
Rips parameter.  Each triangle meets only a uniformly bounded number of
sampled cuts.  Summing over tracks therefore costs only a fixed multiple
of the number of triangles.  The horizontal boundary lengths and the
bounded section lengths give
\begin{equation}\label{gate:corridor}
 \sum_{-d<j<d}l_j\le C(l_{-d}+l_d)+C_m,
 \qquad 1\le d\le m,
\end{equation}
where $C$ is independent of $m$.

Put $P_d=\sum_{|j|\le d}l_j$.  Equation~\eqref{gate:corridor} implies
\[
 P_d\ge (1+C^{-1})P_{d-1}-C_m/C.
\]
With $m$ fixed, iteration shows that the two moving ends at $m$ dominate
$l_0$ by a factor tending exponentially to infinity with $m$, apart from
an additive error depending on $m$.  First choose $m=b$ to make this
factor exceed the desired constant, and then choose the threshold for
$l_0$ to absorb the additive error and the two trimming errors.  Increasing
that threshold to an integer gives $R$.
\end{proof}

Choose $R$ still larger so that every bounded-comparison additive error
and trimming loss used below is small compared with $R$.  Comparisons send
vertices to vertices, and we take integral trimming lengths.  An oriented
arc of length $2R$ has an \emph{expanding event of sign $\varepsilon$}
if its trimmed image under the $\varepsilon b$ comparison has length at
least $3(2R)$.  After such an event, expansion continues in that sign,
with factor at least $100$ at each subsequent $b$-step for which the arc
is carried: immediate inverse comparison is too short to be the expanding
direction in Lemma~\ref{gate:flaring}.  By partitioning a long expanded
arc into length-$2R$ pieces, the triangle inequality and the uniform upper
Lipschitz bound show that arbitrarily many disjoint expanding events occur
far along an infinite expanding trajectory.

The next lemma isolates the cyclic obstruction used in phase uniqueness.
It also supplies the infinite-intersection assertion needed later: a
rational line surviving all finite windows will have a common nonzero
translation. The finite datum in the proof is an actual pair of
continuation windows with a marked edge, not an isomorphism type of
free groups.

\begin{lemma}[Cyclic indices along a trajectory]\label{gate:cyclic-indices}
Assume the expansion package. Fix a bounded sampling step, and let
\(W_j\) be the successive short windows, with sampled edges \(e_j\),
along an interval of a monotone trajectory. The interval may
be finite, one-sided infinite, or bilateral. Put
\[
 H_j=\Stab(W_j),\qquad
 K_j=\Stab\bigl(\operatorname{Hull}(W_j\cup W_{j+1})\bigr).
\]
Let \(C=\langle c_0\rangle\le L\) be a maximal cyclic subgroup of
\(G\). Suppose \(C\cap K_j\) is nontrivial and loxodromic in the
short-window trees for every adjacent pair. Write
\[
 C\cap H_j=\langle c_0^{q_j}\rangle,\qquad
 C\cap K_j=\langle c_0^{m_j}\rangle,
 \qquad q_j,m_j>0.
\]
There are uniform constants \(B,E\) such that both indices
\(m_j/q_j\) and \(m_j/q_{j+1}\) are at most \(B\), and at most
\(E\) of these indices differ from one over the whole interval.
In particular all partial products of their ratios lie between
\(B^{-E}\) and \(B^E\). For an infinite interval some nonidentity
element of \(C\) belongs to every \(K_j\).
\end{lemma}
\begin{proof}
Let \(V_j=\operatorname{Hull}(W_j\cup W_{j+1})\) and set
\(r_j=c_0^{q_j}\). The axis of \(r_j\) in \(R_j\) lies in the
\(K_j\)-core, since a nonzero power of \(r_j\) belongs to \(K_j\).
That core has uniformly finitely many edge orbits for this bounded
window shape. Among successive translates of a marked oriented
axis edge by \(r_j\), two therefore belong to the same \(K_j\)-orbit
after uniformly many steps. An oriented edge has trivial stabilizer,
so the element identifying them is the corresponding power of
\(r_j\). This bounds \([C\cap H_j:C\cap K_j]=m_j/q_j\).
The argument in \(R_{j+1}\) bounds the other index.

Suppose \(m_j/q_j>1\). Then \(r_j\) fixes \(W_j\) but not
\(V_j\): a height-zero element preserving the finite monotone
window setwise fixes it pointwise. Consequently \(V_j\) and
\(r_jV_j\) are distinct actual continuations through \(W_j\),
with their terminal positions at the same height. Their stabilizers
\(K_j\) and \(r_jK_jr_j^{-1}\) both contain \(C\cap K_j\).
By (E4) their intersection is cyclic, and its maximal cyclic
overgroup in \(G\) is exactly this same \(C\).

Mark an oriented regular edge on their common axis in \(R_j\).
There are only finitely many orbits of the resulting actual marked
pairs. Indeed (E2) gives finitely many continuation-window orbits;
their finite core edge quotients give finitely many orbits after
adding an edge marking. Normalize the short window and that edge.
Since its stabilizer is trivial, only finitely many actual lifts of
each continuation remain. Hence only finitely many pairs remain,
with their actual subgroup embeddings retained. The same argument
applies to a nonunit index at the \(W_{j+1}\)-end, recording which
end is marked.

If one marked-pair type occurred at two different heights, its
transporter \(g\in G\) would conjugate the two cyclic intersection
groups just described. Both have maximal cyclic overgroup \(C\)
as an actual subgroup of \(G\), so \(g\) would commensurate \(C\).
But equivariance of height gives
\(h(g)=h(e_{j'})-h(e_j)\ne0\), which is impossible. A station has
only the two adjacent comparison directions. The finite list thus
bounds the total number of nonunit indices on the whole interval,
including when the interval is all of \(\mathbb Z\). This proves
the bound \(E\) and the partial-product assertion.

Except at these finitely many indices,
\(q_j=m_j=q_{j+1}\). Thus only finitely many different integers
\(q_j,m_j\) occur, even on an infinite interval. A common multiple
\(M\) gives \(1\ne c_0^M\in K_j\) for every \(j\).
\end{proof}

\begin{corollary}[Persistence of a rational line]\label{gate:rational-persistence}
Suppose a bilateral trajectory carries, in \(R_0\), the whole axis
of a loxodromic element \(c\in H_0\) through every finite subwindow
containing position zero. Some nonzero power of \(c\) fixes the
entire trajectory pointwise.
\end{corollary}
\begin{proof}
For each finite window, its core contains this axis and has finite
edge quotient. Repetition of a marked axis edge under powers of
\(c\) therefore puts a nonzero power of \(c\) in its stabilizer.
This conclusion initially allows the power to depend on the window.
Equivariant comparisons transfer the axis to any other station as
the axis of that same actual element, after taking a power lying
in the comparison group. Thus the maximal cyclic overgroup \(C\)
of \(c\) meets every adjacent comparison group nontrivially and
loxodromically. Lemma~\ref{gate:cyclic-indices}, applied over all
\(\mathbb Z\), supplies one common nonzero power. The comparison
windows cover the trajectory, so that power fixes it pointwise.
\end{proof}

\begin{lemma}[Phase uniqueness]\label{gate:phase-uniqueness}
There is a uniform depth $N$ such that two trajectories carrying the same
expanding event of length $2R$, and agreeing through depth $N$ in its expanding direction
including the short comparison windows, cannot subsequently diverge in
that direction while still carrying the arc.  One or both trajectories
may be finite.  The same conclusion holds for arcs carried by all finite
subwindows of an infinite trajectory.
\end{lemma}
\begin{proof}
The obstruction to late divergence is a long periodic overlap. We
pull its periods back to the original event. Adjacent primitive
periods can change by a bounded index, but
Lemma~\ref{gate:cyclic-indices} bounds the number of such changes.
Exponential expansion will therefore exceed the total loss in this
pullback.

Suppose there is a first divergence.  At the last available common
$b$-station before it, the expanded arc lies in the cores of two bounded
continuation groups.  The continuation windows need extend only far
enough to exhibit the first divergent edges; the distance from the station
to these edges is bounded.  No comparison beyond the divergent edges is
required.  Their intersection is at most cyclic by (E4).

We record the finite-state overlap argument.  For an oriented edge in the
intersection of two cores, record its orbit in each core.  There are
uniformly finitely many such pairs.  If two edges have the same pair of
records, the two elements carrying one edge to the other coincide, because
an oriented edge has trivial stabilizer.  Their common value belongs to
the intersection of the two groups.  Thus the overlap has uniformly
finitely many edge orbits under that intersection. If the intersection
is trivial, the overlap has uniformly bounded diameter. A nontrivial
cyclic intersection can also be elliptic in a relative tree. In that
case project one of its fixed vertices to the invariant overlap; the
projection is fixed too. A radial segment from this point contains at
most one edge of each cyclic orbit, since the action preserves distance
from the fixed point. The finite edge-orbit bound therefore bounds the
radius of the overlap. Thus a sufficiently long overlap forces the
cyclic intersection to be loxodromic. It then runs along its axis apart
from bounded initial and terminal pieces: an off-axis hanging segment
cannot contain a repeated oriented-edge state. Its period is uniformly
bounded.

Let $r_j$ be the primitive translation of that axis in the short-window
group at the last common station $j$, oriented with the arc.  The arc
contains exponentially many complete periods.  Pull these periods back
through the preceding comparison group.  If sufficiently many remain,
repetition in its finite-edge-quotient core gives
\[
 r_j^a=r_{j-b}^{a'}
 \quad\hbox{as actual group elements},\qquad
 1\le a,a'\le C',
\]
where each side is the least positive power of the corresponding primitive
root lying in the comparison group.  Exact ellipticity ensures that the
common element remains loxodromic on the other side.  Repetition bounds
the indices by the uniform edge-orbit count.  Equivariance of the coarse
inverse on the axis implies that if $p$ periods remain on the later side,
there are at least $(a'/a)p-O(1)$ periods on the earlier side.  Removing
partial blocks and end margins costs only a uniformly bounded number of
periods.

Every equality above identifies powers as actual elements of \(G\).
All the roots reached in the pullback therefore lie in one maximal
cyclic subgroup \(C\le L\). Lemma~\ref{gate:cyclic-indices} applies
on every interval already reached, with bounds independent of its
length. In particular the product of the ratios \(a'/a\), on any
part of the pullback and in either direction, lies between
\(B^{-E}\) and \(B^E\).

Let \(t\) be the number of backward steps to the initial event and
\(p_t\) the number of terminal periods. The recurrence just proved
gives, after \(t\) steps,
\[
 p_0\ge B^{-E}p_t-C_0B^E t
\]
for a uniform trimming constant \(C_0\); the same lower bound with
fewer losses holds at every intermediate station. The terminal
period has uniformly bounded length, while the event has expanded
by at least \(100^t\). Hence \(p_t\ge c_1 100^t-C_1\), with
\(c_1>0\) uniform. For sufficiently large \(t\), all intermediate
period counts exceed the finite repetition threshold required to
take the next backward step, and the displayed lower bound exceeds
\(2R\). But each initial period has length at least one, whereas
the initial event has length \(2R\). This is the contradiction.
Choosing the common depth \(N\) to allow this many steps proves the
lemma. All terminal constants came from the bounded windows
exhibiting the first divergence, so the proof is uniform for finite
trajectories and for finite-window carrying of infinite trajectories.
\end{proof}

\begin{lemma}[Changing the sampling offset]\label{gate:offset}
For an ordinary bilateral trajectory, an expanding event of a given sign
produces events of that sign at every sufficiently distant row in that
direction, not merely on its stride-$b$ subsequence.
\end{lemma}
\begin{proof}
There are only finitely many offsets of the stride-$b$ sampling.  Fix one.
If at a sufficiently distant row on this offset there were no events of
the given sign on the transferred lines, partitioning long arcs into
length-$2R$ pieces bounds their transfer in that sign.  Lemma~\ref{gate:flaring}
then forces expansion by at least $100$ in the opposite sign, and this
expansion persists.

Start with a long arc already expanding on the original subsequence,
transport it forward $m$ steps and then across the bounded offset, and
transport it backward $m$ steps on the offset subsequence.  The resulting
trimmed arc lies, up to bounded error, inside the direct bounded-offset
transfer of the starting arc.  This follows from
Lemma~\ref{gate:comparisons} on the common finite diagram, and stability
of lines under the comparisons.  Choose $m$ large enough that the two
opposed exponential expansions exceed the upper Lipschitz constants of
the bounded-offset comparisons.  With this $m$ fixed, move the initial
row far enough along the original expanding trajectory to dominate every
additive error for that finite diagram.  The length comparison is then
impossible.  Taking the maximum over the finitely many offsets proves
the assertion.
\end{proof}

\section{Actual trajectory clusters}\label{gate:trajectory-clusters}

Outside the acylindrical case, we use expansion to construct a
$G$-invariant system of periodic quasiconvex groups, finite modulo $G$,
whose height kernels contain all sufficiently long noncyclic path
stabilizers.
We work in the coarse actual tree $T_0$.  The ordinary expansion
package follows from the actual splitting and path-inertia statements of
Section~\ref{sec:03-rows}.  In particular the total ordinary space is
quasi-isometric to $G$, the horizontal trees are locally finite free Cayley
trees, and Proposition~\ref{row:path-inertia} supplies joint homology injection
and its subgroup-wise version.  The sum of reduced ranks for paths of a
fixed shape is at most $k-1$.  The induced-coefficient content of that
lemma, not merely this numerical bound, will be used below.

There are three steps. We first obtain periodic end stabilizers from
branching surviving subtrees. We then connect positive and negative
ends when they support a common bilateral trajectory. Finally, we
show that each connected component of this switch graph is a tree
and that its stabilizer captures every sufficiently long path group
seen by that component.

\subsection{Surviving cores and periodic ends}

For an infinite monotone path all of whose relevant finite windows have
noncyclic stabilizer, intersect their decreasing cores in a fixed
horizontal tree.  Call the intersection the \emph{surviving subtree}.
We only consider an event or a branch when it belongs to this intersection.
Local finiteness implies that any nonempty surviving subtree is a union
of full lines: extending a finite surviving arc in each finite-window core
and taking subsequences gives extensions in both directions.  The same
argument preserves a specified surviving arc.  For a bilateral path,
Lemma~\ref{gate:comparisons} transfers the endspace of its surviving
subtree between rows.

This is a geometric intersection of subtrees, not yet the core of
the intersection of their groups. A surviving line need not initially
come with a common nontrivial translation. A trajectory or tail is
called \emph{branching} when its surviving subtree has a branch
vertex, equivalently three distinct ends. Comparisons make this
condition independent of the sampled row. The next approximation
lemma supplies the group-theoretic control needed alongside these
geometric data.

\begin{lemma}[Decreasing core graphs]\label{gate:decreasing-cores}
Let \(K_n\) be a decreasing sequence of finitely generated subgroups
of one free row, with uniformly bounded ranks, and put
\(M=\bigcap_n K_n\). Then \(M\) is a finite-rank free group. Its
core graph embeds in that of \(K_n\) for all sufficiently large
\(n\), and its first field homology injects there.
\end{lemma}
\begin{proof}
Choose a finite subgraph of the core graph of \(M\), and choose
actual horizontal lifts of its vertices and edges. An identification
of two chosen vertices in a \(K_n\)-quotient has one possible
transporter, because a horizontal vertex has trivial stabilizer.
If it does not belong to \(M\), it fails to belong to some \(K_n\)
and to every later group. There are only finitely many candidate
identifications, so the chosen finite subgraph eventually embeds.
An infinite-rank core graph would have finite subgraphs of arbitrarily
large rank, contradicting the uniform rank bound on \(K_n\). Thus
\(M\) has finite rank. Apply the same argument to its entire finite
core graph. An embedded connected subgraph gives a free factor,
after choosing a connecting path to the basepoint, and consequently
an injection on first field homology.
\end{proof}

\begin{lemma}[Branch states and tail groups]\label{gate:tail-groups}
There are finitely many orbits of branching bilateral trajectories and of
branching one-sided tails.  Every branching bilateral trajectory is
periodic under a nonzero-height element; every branching one-sided tail
is eventually periodic.

For a monotone ray from an edge $e$ to an end $z$, let
\[
 M_{z,e}=\Stab([e,z))
\]
be its pointwise stabilizer.  This is a finite-rank free group.  Its
homology injects in sufficiently long finite-prefix homologies, and the
finite-path restriction rule also holds on restricting within a later
tail group, with the same common-level and orbit conventions.  If some
$M_{z,e}$ is noncyclic, let $Y_z$ be the union
of all monotone rays to $z$ having noncyclic stabilizer, and let
$D_z=\Stab_G(z)$.  Then:
\begin{enumerate}[label=(\roman*)]
\item $Y_z$ is connected, $D_z$ has a height translator with axis in $Y_z$,
and tail ranks along such an axis are constant;
\item $D_z\cap L$ is transitive on the $Y_z$-edges of each fixed level;
\item $D_z$ is quasiconvex in $G$ and virtually an ascending mapping torus
of a finite-rank free group; in particular it is virtually compact special.
\end{enumerate}
\end{lemma}
\begin{proof}
We first produce translators using a finite set of trajectory orbits.
The core-graph lemma then gives the tail stabilizers and their
restriction maps. These inputs establish transitivity along levels
and finally ambient quasiconvexity.

In a noncyclic finite-rank free core, the number of branch-vertex orbits is
at most twice the reduced rank.  At any finite depth, count pairs consisting
of a path and a branch vertex of its core, modulo the horizontal free
vertex action.  Joint path inertia bounds this count uniformly by
$2(k-1)$.  A specified horizontal vertex has trivial stabilizer, so marking
it removes all ambiguity in the transporter. Normalize the sampled edge
to a fixed representative edge at the chosen level and a surviving
branch vertex to one horizontal vertex.
More than \(2(k-1)\) distinct branching infinite paths in these
coordinates would already be distinct at some common finite depth.
Thus there are finitely many \(L\)-orbits of branching bilateral
trajectories with a sampled edge at this level, and finitely many
\(L\)-orbits of branching tails starting there. These unmarked
trajectory orbits are the states; branch markings only bound their
number and are not part of a proposed canonical shift of vertices.

Choose a height-one element \(t\). For a positive trajectory, move
the sampled position to its next edge and translate the whole picture
by \(t^{-1}\) to restore its height. On bilateral states this gives
a bijection, with inverse obtained from the preceding edge and
\(t\). On positive tail states first delete the initial segment,
then translate by \(t^{-1}\); this is a map of the finite state set.
Both operations are well defined on \(L\)-orbits since \(L\) is
normal, and preserve branching because comparisons preserve the
three ends of a surviving tripod. A repeated bilateral state gives
an element of nonzero height preserving the same actual trajectory.
A repeated tail state gives an element taking a later tail to a
proper subray of itself, hence translating its eventual axis.
Negative tails use the reversed operation. This proves periodicity
and eventual periodicity without choosing a transported branch vertex.

The tail group \(M_{z,e}\) is the intersection of its decreasing
finite-prefix groups. Their ranks are bounded by path inertia, so
Lemma~\ref{gate:decreasing-cores} proves finite rank and the asserted
homology injection into sufficiently long prefixes.
For a finite family in a prescribed backward restriction sum, choose
these depths simultaneously and compose with
Proposition~\ref{row:path-inertia}.  When the restriction takes place within a
later tail group $M_{z,f}$, use the subgroup-wise version of that lemma:
intersection with the finite segment from an earlier edge to $f$ is
exactly the earlier tail group.  This proves the precise inherited
restriction rule used below.

Rays to a common end coalesce; their union $Y_z$ is consequently connected.
A noncyclic tail group gives a branch in the surviving subtree.  Eventual
periodicity of that tail produces a height translator in $D_z$.  Its axis
has a tail in $Y_z$, and invariance under its powers puts the whole axis
in $Y_z$.  Tail groups increase as the starting edge moves towards $z$.
The inherited inertia gives injective first-homology maps for these
inclusions.  Periodicity makes their dimensions repeat, so those
dimensions, and hence their ranks, are constant along the axis.

Compare an arbitrary $Y_z$-edge at a given level with the axis edge at
that level.  Their rays meet at a later axis edge $f$.  Inside $M_{z,f}$,
the axis predecessor already contributes its full first-homology dimension.
The joint restriction rule within $M_{z,f}$ therefore excludes a second
noncyclic predecessor orbit.  A transporter realizing the common orbit
belongs to $M_{z,f}\le D_z\cap L$.  This proves (ii).

A height-zero element preserving $z$ fixes a tail of every ray to $z$.
Consequently $D_z\cap L$ is the increasing union of tail stabilizers along
the axis.  After passing to a finite-index subgroup to select a height
period, a translator together with one finite-rank tail group gives the
ascending mapping-torus presentation.  At a vertex of $Y_z$, the
stabilizer equals the stabilizer of its parent edge towards $z$.  It is
quasiconvex in the vertex row by the preceding finite-rank assertion and
Lemma~\ref{row:row-finiteness}.  An edge leaving $Y_z$ has at most cyclic
intersection with $D_z$, since a noncyclic such intersection would create
a second outgoing continuation, contrary to the reversal restriction.
The action on $Y_z$ has finite quotient by (ii) and periodicity.
Lemma~\ref{row:boundary-gates} now gives ambient quasiconvexity of $D_z$.
It is therefore hyperbolic, and Theorem~\ref{asc:vcs} makes it virtually
compact special.
\end{proof}

\subsection{Active colors and switches}

Color an expanding event on an actual bilateral trajectory by its expanding
end and its sign.  Lemma~\ref{gate:phase-uniqueness} says that sufficiently
large fixed-window data determine this color uniquely.  Such colors are
called \emph{active}.  A \emph{switch trajectory} is a bilateral trajectory
whose surviving subtree supports events of both signs at some row.

\begin{lemma}[Periodic switches and noncyclic active tails]\label{gate:switches}
Every switch trajectory is periodic.  For an active color $z$ occurring
at an edge $e$, the tail group $M_{z,e}$ is noncyclic.  There are finitely
many orbits of active colors and of switch trajectories.  A switch
trajectory has both signs at every row.
\end{lemma}
\begin{proof}
Every vertex of a surviving subtree lies on a full line and admits a
centered event in at least one sign, by flaring.  If both signs occur at
a row, connect their centers in the surviving tree.  There are opposite
centered events at distance at most one: along the connecting vertex path
the available sign changes, or a vertex admits both signs.  The fixed-window
data of such a nearby pair have finitely many orbits, and phase uniqueness
determines both ends from these data.

Suppose first that a switch has both signs at infinitely many rows of a
stride-$b$ subsequence.  A repeated nearby-pair state determines the same
two ends and hence the same bilateral path.  The change of height gives
a translator, proving periodicity.  Otherwise, persistence of its initial
positive and negative events gives purely positive rows far in the
positive direction and purely negative rows far in the negative direction.

At a pure row, a surviving line lies in a bounded neighborhood of the
orbit of the corresponding tail group.  Indeed, its centered events have
only finitely many marked types.  An element transporting one occurrence
of a type to another fixes the marked row and preserves the uniquely
determined expanding end.  It therefore belongs to the tail group.
Choose one representative for each occurring type.  Every vertex of the
line lies within a fixed distance of a translate of one of those finitely
many representatives.  This proves the assertion and puts the two endpoints
of the line in the limit set of the tail group.

This limit-set conclusion transfers back to an earlier row.  Intersect
the later tail group with the finite connecting-path group.  The result
is the earlier tail group, expressed in the later frame.  Both groups
are finitely generated subgroups of a free row, hence quasiconvex, and
their limit-set intersection is exactly the limit set of their intersection.
The equivariant comparison transfers this intersection of endspaces back
to the earlier row.  Apply the argument on both sides of the switch.
The endpoints of a central surviving line then lie in the limit set of
the bilateral pointwise stabilizer.

If the surviving subtree is not already branching, it is a line; the last
conclusion makes it a rational cyclic axis.  A nontrivial translation of
that axis fixes every window.  Every finite-window group is noncyclic,
so its core branches somewhere on this axis.  Translate the branch point
into a fixed translation period.  Local finiteness and nesting leave a
branch direction present at every depth.  This gives a surviving branch,
a contradiction.  The trajectory is thus branching and periodic by
Lemma~\ref{gate:tail-groups}.

Now consider an active event on a periodic trajectory.  Expansion produces
arbitrarily many disjoint events of the same sign.  Translate them back
by periods to one central row.  Repeated marked event types at different
horizontal positions give a nontrivial element of $M_{z,e}$.  This group
cannot be cyclic: its tail groups are nested and commensurable if cyclic,
and the period translator would commensurate a nontrivial horizontal
cyclic subgroup.

For completeness, a nonperiodic single-sign event requires a separate
argument.  Far along its expanding stride-$b$ subsequence the rows are
pure in that sign; otherwise the nearby-pair argument would make the
trajectory periodic.  The pure-row limit-set argument and transfer back
put the endpoints of every surviving line in the limit set of $M_{z,e}$,
so this group is nontrivial.  If it were cyclic, the line would be rational.
That rational line belongs to every finite-window core by the
definition of the surviving subtree. Corollary~\ref{gate:rational-persistence}
therefore supplies one nontrivial translation fixing the whole
bilateral trajectory. In particular this conclusion uses the index
bound on both directions together, rather than choosing a different
power for each finite window.
The preceding branch-in-a-period argument again produces a surviving
branch, and hence periodicity, a contradiction.  Therefore every active
tail group is noncyclic.

A nearby pair determines a switch, so there are finitely many switch
orbits.  A single fixed-window event determines its color, giving finitely
many color orbits.  Finally Lemma~\ref{gate:offset} supplies both signs at
all sufficiently far appropriate rows on a periodic switch.  Translation
by periods brings those signs to any prescribed row.
\end{proof}

Form the actual bipartite \emph{switch graph} $\mathcal C$: its vertices
are active positive and negative ends, and each switch trajectory joins
its two ends.  Retain isolated active vertices.  A pair of ends determines
at most one bilateral path in the tree, so there are no parallel edges.
The graph has finitely many vertex and edge orbits by
Lemma~\ref{gate:switches}.

\begin{lemma}[The switch graph seen by a long window]\label{gate:local-switch-tree}
There is a uniform depth such that the following holds for every centered
monotone window $W$ whose two endpoints are at least that far from its
sampling center and whose pointwise stabilizer is noncyclic.  The colors
realized in the central core of $\Stab(W)$ are exactly the active ends $z$
with $W\subset Y_z$.  These vertices, together with the switch trajectories
traversing $W$, form a nonempty connected tree, denoted $\mathcal C(W)$.
If it is an isolated vertex locally, that vertex is isolated in
$\mathcal C$.
\end{lemma}
\begin{proof}
First justify the passage from a finite window to bilateral witnesses.
A local state consists of a centered event, or a pair of events with centers
at distance at most one, together with a fixed-size path neighborhood
large enough for phase uniqueness.  There are finitely many normalized
states.  More generally, fixing an actual regular horizontal marking,
there are finitely many carried extensions of any prescribed finite
depth: the relevant paths have finitely many orbits and their cores have
finite edge quotients, while the marking has trivial stabilizer.  Thus
the extension sets form a finitely branching inverse system.

If a given state occurs in arbitrarily deep window cores, choose compatible
extensions successively.  The resulting bilateral path carries its
specified event or pair in every finite subwindow.  Local finiteness of
the horizontal tree preserves the full lines carrying those arcs.  States
that do not persist in this way have a maximum possible depth.  Taking
a maximum over the finitely many short states gives a uniform depth beyond
which every state in question has a bilateral witness.  The witness is
initially required to agree only on the fixed-size neighborhood recorded
in the state.  Lemma~\ref{gate:phase-uniqueness}, applied to the witness
and the finite window, then forces its expanding ray to agree throughout
the corresponding side of $W$.  This last step is what permits use of a
fixed-size state for an arbitrarily long $W$.

Every vertex in the finite-window core admits a centered event.  Assign
its expanding color using a bilateral witness.  Phase uniqueness makes
the assignment unambiguous and equivariant under $\Stab(W)$.  Adjacent
same-sign labels agree: use a common bilateral witness for the pair.
Adjacent opposite-sign labels give a switch through $W$.  Connectedness
of the core therefore connects all recorded colors in the switch graph.
A recorded color either participates in a switch or is the unique color
and is preserved by the noncyclic window group.  In the switch case,
periodicity and Lemma~\ref{gate:switches} imply $W\subset Y_z$.  In the
unique-color case, the ray to that end extends $W$ monotonically; the
window group fixes that end and $W$, and hence its entire continuing tail.
The tail stabilizer is noncyclic, so again $W\subset Y_z$.

Conversely, suppose $W\subset Y_z$.  Take events far along an expanding
trajectory for $z$, and use periods and Lemma~\ref{gate:offset} to bring
them to the required levels.  By the transitivity in
Lemma~\ref{gate:tail-groups}, match the starting edge on the side away
from $z$.  The path towards $z$ from that edge is unique.  Taking the
original event far enough along its expanding part ensures that its
comparison subwindow lies on this ray in $Y_z$; it therefore realizes
the desired event in the central core of $W$.  The same argument applied
to switch paths shows that a locally isolated occurrence has no global
switch incidence.

It remains to exclude a cycle.  At the center, use the convex cores of
the tail groups $M_{z,e}$ for the colors just obtained, with one fixed
thickening.  Nearby paired events are uniformly close to the corresponding
tail cores.  To verify uniformity, normalize their marked fixed-window
data; there are finitely many types, each determining the actual end,
and the tail group acts coboundedly on its core.  Every switch through
$W$ has such nearby paired events by periodicity and
Lemma~\ref{gate:offset}.  Thus every switch edge gives an intersection
of the two thickened cores.

For sufficiently deep $W$, the thickenings of distinct same-sign cores
are disjoint.  Otherwise normalize a pair at bounded distance near one
horizontal origin.  There are only finitely many such actual lifts: use
finite color orbits, edge transitivity on $Y_z$, and coboundedness of each
tail core.  Increasing depths then give arbitrarily long agreement of
the corresponding rays, forcing equality of their ends by phase
uniqueness.  This contradicts distinctness.  Finally, the intersection
graph of subtrees of a tree is chordal, by convexity and the Helly
property.  With no same-sign intersections it is bipartite, hence a
forest.  The switch graph under consideration is a subgraph of this
intersection graph; extra opposite-sign intersections cause no problem.
It is connected by the preceding argument, and therefore is a tree.
\end{proof}

\subsection{Splitting marks and simultaneous free factors}

We next pass from individual periodic ends to the full stabilizer of
a connected component of switches. To see this group through finite
windows, temporarily distinguish incidences at an end according to
where their rays cross a chosen level. Widely separated levels cut
the component into the local trees already constructed. Moving the
levels towards the respective ends joins those pieces. We will show
that these joins neither create circuits nor increase the rank of
their stabilizers. This is where the actual homology injections,
rather than only the rank bound, enter.

Fix a nonisolated component of $\mathcal C$, let $A$ be its full stabilizer,
and put $S=A\cap L$.  For integers $j,l$, split each positive vertex by
also remembering the level-$j$ edge of each incident trajectory, and split
each negative vertex by remembering its level-$l$ edge.  Incidences with
the same end and the same additional mark remain together.  Denote the
resulting graph by $\mathcal C_{j,l}$, retaining its $S$-action.  Increasing
$j$ and decreasing $l$ coarsens the graph: marks move towards the common
end and eventually coincide.  All these actions preserve the bipartition
and have no inversions.

For $j$ sufficiently smaller than $l$, its components are precisely the
local trees $\mathcal C(W)$ for the paths across $[j,l]$ just considered.
Their stabilizers in $S$ are exactly
\[
 K(W)=\Stab(W).
\]
Indeed the two marks determine the segment $W$, and equivariance gives
both inclusions.  The split vertex stabilizers are $M_{z,e}$ and the split
edge stabilizers are
\[
 M_\tau=\Stab(\tau),
\]
the pointwise stabilizers of the bilateral switch trajectories.  The same
finite-core approximation used for tails shows that $M_\tau$ is finitely
generated free.  It is noncyclic or trivial: a nontrivial cyclic $M_\tau$
would be normalized by a period translator of $\tau$, contrary to cyclic
noncommensuration.

\begin{lemma}[Orbit control under coarsening]\label{gate:mark-orbits}
The components of each $\mathcal C_{j,l}$ form one $S$-orbit.  Their
quotient vertex, edge, and incidence counts are finite and do not change
as $j$ increases and $l$ decreases, including after the two levels cross.
The corresponding vertex and edge ranks do not change either.
\end{lemma}
\begin{proof}
Consider two incident switch trajectories sharing an end.  Orient their
period translators towards that end, so their heights are negative when
the common end is negative, and take positive powers having equal heights.
Far enough
towards the shared end, any prescribed finite subwindows of the two
trajectories lie on their common ray.  Move the first subwindow there
using its period translator and move back using the other translator.
Both translators stabilize the original component because they preserve
an edge in it.  Their composite therefore belongs to $A$; it has height
zero, hence belongs to $S$, and identifies the relevant split components.
Connectivity of the
original component propagates this comparison along any finite path of
switch incidences.  Hence all split components lie in one $S$-orbit.

For a fixed end, $D_z\cap L$ is transitive on the eligible marked edges
at a level by Lemma~\ref{gate:tail-groups}; their stabilizer ranks are
constant at all levels. More explicitly, an \(A\)-orbit of ends
contributes
\[
 [h(A):h(D_z)]
\]
\(S\)-orbits of marked ends, independently of the marking level.
An \(A\)-orbit of switch trajectories contributes
\( [h(A):h(\Stab_A(\tau))]\) \(S\)-orbits of split edges.
The marks on a switch are its uniquely determined edges at the two
chosen levels, and its pointwise group \(M_\tau\) does not change.
Each split edge has exactly one incidence in each bipartition, so
the same edge list also controls the two incidence lists. All these
indices are finite because the stabilizers have height translators.

There are finitely many \(A\)-orbits: if a group element carries
one end or switch in this component to another, it carries their
entire connected components to one another and hence belongs to
\(A\). Thus the finite global orbit lists restrict to finite
\(A\)-lists. Because the split components form one \(S\)-orbit,
the lists just counted are exactly the quotient lists in one
component under its stabilizer. Their counts and ranks are therefore
independent of the levels.
They persist under coarsening, and every resulting component comes from
components at a sufficiently separated pair of levels.
\end{proof}

\begin{lemma}[Simultaneous incident free factors]\label{gate:incident-factors}
At every vertex of every split action, the nontrivial incident edge groups,
indexed by incidence-orbit representatives, are simultaneous free factors
of its vertex group, up to separately chosen inner conjugations.
\end{lemma}
\begin{proof}
Fix a marked end $(z,e)$ and write $M=M_{z,e}$.  Start a long path at $e$
towards $z$, and compare at a center sufficiently far along it and far
from both endpoints.  By periodicity and the finite number of switch
orbits, the core graphs of the nontrivial $M_\tau$ have uniformly
bounded sizes at these centers.

In the core graph of the finite-path group $K$, include the core graphs for
\emph{all} switch lines through that path.  They eventually embed jointly
and disjointly modulo their $K$-orbits.  Here is the identification test.
An alleged overlap of chosen actual horizontal vertices has a unique
transporter.  Normalize at that vertex.  Finite switch orbits and the
uniform core-size bound leave only finitely many normalized switch lifts
nearby.  Distinct such bilateral paths have only a bounded common
centered subpath.  Once the comparison path is longer, any persistent
identification must preserve the entire switch line.  An element of $K$
has height zero, so if it preserves that line it fixes it pointwise.
Consequently its line stabilizer is exactly $M_\tau$.  These are precisely
the identifications already made in the $M_\tau$ core graph.  This
proves both embedding and disjointness.  The same finite test, applied
to the whole finite collection of core representatives, is uniform.

Disjoint embedded connected core subgraphs in a free graph give
simultaneous free factors of its fundamental group, after choosing
connecting paths.  Therefore the indicated groups are simultaneous free
factors of $K$.  Since $M\le K$, apply the Kurosh subgroup theorem to
this free-product decomposition.  Within a $K$-orbit of switch lines,
the desired incidences modulo $M$ correspond to distinct double cosets
\[
 M\backslash K/M_\tau.
\]
Their pointwise groups are already contained in $M$, so the Kurosh
intersection is the entire conjugate edge group.  This yields exactly
the requested simultaneous free factors in $M$, with the independent
conjugations arising from its choices of connecting paths.
\end{proof}

\begin{lemma}[A homology obstruction to a proper free quotient]\label{gate:free-quotient}
Let $F_r$ map onto a nontrivial subgroup $B\le L$, and use $D_L$
coefficients.  If this map has nontrivial kernel, then
\[
 \dim_{D_L}H_1(B;D_L)\le r-2.
\]
In particular, it cannot surject onto a group containing a subgroup whose
ambient homology image has dimension $r-1$.
\end{lemma}
\begin{proof}
Choose a shortest nontrivial freely reduced word in the kernel.  Its
proper prefix vertices in the quotient are distinct; a repeated one
would give a shorter nontrivial kernel word.  Its closed labelled path
has no cyclic backtrack, since cyclic cancellation would again shorten
the kernel word.  At least one evaluated Fox derivative is nonzero.
More explicitly, cancellation of a contribution to a generator derivative
would require an opposite traversal of the same oriented labelled edge.
This repeats prefix vertices unless it is the forbidden backtrack.
Thus this relator supplies a nonzero 2-boundary over the injective group
ring inside $D_L$.  The 1-boundary for the $r$ generators has rank one,
since $B$ is nontrivial.  Consequently its kernel has dimension $r-1$,
and the nonzero 2-boundary lowers first homology by at least one.  Adding
any remaining relations can only lower that dimension further.
\end{proof}

\begin{proposition}[The component trees]\label{gate:component-tree}
Every component of the actual switch graph is a tree.  For a nonisolated
component with stabilizer $A$ and $S=A\cap L$, there is an integer $r\ge2$
such that $S$ is a nested union of rank-$r$ free groups and
\[
 \dim_{D_L}H_1(S;D_L)=r-1.
\]
These approximants and $S$ have injective ambient first-homology maps.
Moreover $A$ is finitely generated, quasiconvex, virtually compact special,
and field-acyclic.
\end{proposition}
\begin{proof}
There are two steps to excluding circuits. Simultaneous incident free
factors make the abstract incidence graph of groups free of a fixed
rank \(r\). Its natural quotient onto a component stabilizer cannot
be proper, because an old path subgroup already contributes
\(r-1\) independent ambient homology classes. We apply this first
to finite coarsenings, then exhaust the unsplit component.

For separated marks, a split component is the tree $\mathcal C(W)$ and
its stabilizer is the noncyclic free group $K(W)$.  Put $r=\rk K(W)\ge2$.
For arbitrary coarsened marks, form the abstract incidence graph of groups
for a component and its stabilizer.  Lemma~\ref{gate:incident-factors}
makes its fundamental group free: eliminate each glued free factor on
one side of an edge, retaining the remaining free factors and the usual
stable letters.  Its rank is constant under coarsening, because the
finite quotient incidence counts and all vertex and edge ranks are
constant by Lemma~\ref{gate:mark-orbits}.  At separated marks this rank
is $r$.

For a group acting without inversions on a connected graph, the incidence
graph-of-groups fundamental group surjects onto the acting group, with
kernel the ordinary fundamental group of the acted-on graph.  This is
seen by lifting the action to its universal covering tree.  Apply this
to a coarsened component.  Its stabilizer contains an old $K(W)$, whose
homology injects into $H_1(L;D_L)$ with dimension $r-1$ by path inertia.
Lemma~\ref{gate:free-quotient} excludes a proper quotient of the abstract
rank-$r$ free group.  Hence the surjection is an isomorphism and the
coarsened component is itself a tree.

Every finite circuit in the unsplit component would lift to a sufficiently
coarsened split graph: move its finitely many shared-end marks far enough
to make the necessary marks coincide.  Since no split graph has a circuit,
the unsplit component is a tree.  Track split components containing one
chosen switch line.  Their stabilizers increase.  Every $s\in S$ sends
that line to a line connected to it by a finite path in the original
component; the finitely many required marks eventually coalesce.
Therefore $s$ belongs to one tracked stabilizer, proving that their union
is $S$.  Each is free of rank $r$, contains an old $K(W)$ with full ambient
homology dimension, and hence has injective ambient homology of dimension
$r-1$.  The maps between them, and their direct limit, preserve that
injection.

The quotient of the component tree by $A$ is finite.  Its vertex groups
are the end groups $D_z$, already finitely generated and quasiconvex.
An edge stabilizer is the setwise stabilizer of a switch trajectory and
fits into
\[
 1\longrightarrow M_\tau\longrightarrow \Stab_A(\tau)
 \xrightarrow{h}h(\Stab_A(\tau))\longrightarrow1.
\]
Its height image is nonzero by periodicity, and $M_\tau$ is finitely
generated, including the possible trivial case.  Thus edge stabilizers
are finitely generated, and so is $A$.

The period rule, Proposition~\ref{row:period-rule}, now applies to $A$: its
height image is nonzero and $0<\beta(S)=r-1<\infty$.  On the refined
actual tree it supplies the finite-quotient hull required by
Lemma~\ref{row:boundary-gates}, with quasiconvex row intersections and
at most cyclic boundary intersections.  Hence $A$ is quasiconvex in $G$.
The switch edge groups are intersections of the two quasiconvex end groups,
so they are quasiconvex in $G$ and in $A$.  The quasiconvex hierarchy theorem
applied to this finite graph-of-groups decomposition, with its virtually
compact special vertex groups, makes $A$ virtually compact special.

Finally, every end vertex group is a finite extension of an ascending
free-group mapping torus.  Every switch edge group has a finite-rank
bilateral kernel and a nonzero height translator; the kernel is trivial
or noncyclic.  These vertex and edge groups are field-acyclic.  One can
compute this directly by the height homology sequence: over the skew
Laurent division field, translation minus identity on the finite-dimensional
kernel homology is invertible, as is seen from its leading Laurent degree.
The same computation covers a cyclic edge group with trivial kernel.
Finite-index passage preserves the vanishing.  Mayer--Vietoris on the
switch-tree splitting now gives field acyclicity of $A$.
\end{proof}

\subsection{The initial captured system}

\begin{proposition}[Capture by periodic quasiconvex groups]\label{gate:capture}
Unless $G$ is already virtually compact special by the acylindrical case,
there is an invariant system of actual objects, finite modulo $G$, with
full stabilizers $A$ and kernels $S=A\cap L$, having these properties:
\begin{enumerate}[label=(\roman*)]
\item $h(A)\ne0$; each $A$ is quasiconvex, virtually compact special,
and field-acyclic.
\item Each $S$ is a nested union of finite-rank free approximants of one
constant rank at least two.  Its field first-homology dimension is positive
and finite, and its ambient homology detects every noncyclic subgroup.
\item For representatives of the actual object orbits under $L$, the maps
\[
 \bigoplus H_1(S;D_L)\longrightarrow H_1(L;D_L)
\]
are jointly injective.  Distinct actual object stabilizers intersect in
at most cyclic groups.
\item There is a uniform length such that the pointwise stabilizer of
every longer noncyclically stabilized actual path, in $T_0$ or in the
refined tree $T$, is wholly contained in one of these kernels $S$.
\end{enumerate}
\end{proposition}
\begin{proof}
Use the components of $\mathcal C$ as objects, taking $D_z$ for an isolated
color.  Nonisolated components have the required finiteness, quasiconvexity,
specialness, and acyclicity by Proposition~\ref{gate:component-tree}.
For an isolated color the increasing tail groups along a translating
axis have constant positive reduced rank and exhaust $D_z\cap L$.
On a sufficiently long path towards that end, the path group preserves
the unique color by Lemma~\ref{gate:local-switch-tree}.  It fixes the
starting edge and the end, hence the whole ray, so it equals the tail
group at the starting edge.  This proves the corresponding approximant
and homology assertions in the isolated case.  Field acyclicity follows
from the same ascending-torus calculation.

To obtain a simultaneous injection, choose finitely many free
approximants from distinct \(L\)-orbits of actual objects. In each
nonisolated component choose an older separated-window subgroup
\(K(W)\) inside the selected coarsened approximant. Both are free
of the same rank, and \(K(W)\) already has full ambient homology
dimension. Its map to the approximant is therefore an isomorphism
on first homology after extending to \(D_L\). The separating marks
can be taken arbitrarily far apart, so make these choices for all
selected components at one common absolute level range. For an
isolated end, use an earlier tail approximant at the required
starting level; constant tail rank gives the same homology image,
and a sufficiently long window realizes that tail group as above.

The resulting windows belong to distinct \(L\)-orbits: their
equivariantly determined local switch trees lie in the selected
distinct component orbits. Proposition~\ref{row:path-inertia} now
injects their homologies jointly into \(H_1(L;D_L)\). Replacing each
window group by its approximant preserves its ambient image and
dimension. Passing to increasing unions, and testing any proposed
relation in a finite collection of approximants, proves (iii)'s
joint injection. Here the indexing is by actual object orbits,
before any possible identification of their abstract subgroup kernels. Every
noncyclic subgroup of one kernel contains a noncommensurable pair in a
single free approximant.  Such a pair is detected in the approximant's
field homology, and the injected ambient map preserves that detection.
This proves (ii).

We spell out the same-orbit intersection issue.  Different $L$-orbit
types cannot have a noncyclic common subgroup: its detected cycle would
contradict their joint injection.  For an $L$-translate of the same type,
write $S=\bigcup_n F_n$ for the free approximants of rank $r$.  Put
$E=D_S\otimes_{\mathbb Q S}\mathbb Q L$.  Coset independence embeds $E$
in $D_L$.  A nonzero $E$-valued cycle for the free group $F_n$ remains a
nonzero class over $D_L$, because its free resolution has no 2-chains.
The ambient injection for $F_n$ then prevents its image from being an
$E$-valued boundary in $L$.  Consequently
$H_1(F_n;E)\to H_1(L;E)$ is injective.  Decompose the induced module into
double cosets and apply Shapiro's lemma; this gives
\[
 \bigoplus_{SgF_n\in S\backslash L/F_n}
 H_1(S\cap gF_ng^{-1};D_S)\ \hookrightarrow\ H_1(S;D_S).
\]
This is the induced-coefficient restriction behind the slot rule;
a dimension estimate alone would not suffice.  It applies to the
coarsened free approximants because their ambient homology maps have
already been proved injective.  The identity double coset
already has dimension $r-1=\beta(S)$, so every other summand has dimension
zero.  Its group is a subgroup of the finite-rank free group $gF_ng^{-1}$,
and therefore is cyclic or trivial.  For $g\notin S$, taking the increasing
union gives $S\cap gSg^{-1}$ locally cyclic.  A locally cyclic subgroup
of a torsion-free hyperbolic group is cyclic: any nontrivial element
puts the whole subgroup in its cyclic centralizer.  This also shows
$N_L(S)=S$, since $S$ is noncyclic.  In particular equal kernels do not
create distinct objects hidden by a change of representatives.

Now let $C$ be the intersection of two distinct full stabilizers $A$.
Its subgroup $C\cap L$ is at most cyclic by what was just proved.  If
$h(C)=0$, this proves the desired assertion.  If $h(C)\ne0$ and
$C\cap L$ were nontrivial, a nonzero-height element of $C$ would normalize
a nontrivial height-zero cyclic group, contradicting cyclic
noncommensuration.  Thus $C\cap L=1$ in this case, and $C$ embeds in
$\mathbb Z$, so is again at most cyclic.

Finally, a sufficiently long noncyclically stabilized coarse path has a
deep centered window $W$.  Its nonempty connected label graph
$\mathcal C(W)$ lies in a single original component.  The path stabilizer
preserves this graph, hence that component, and has height zero because
it fixes the path.  It is therefore contained in the component kernel.
The same argument includes an isolated color.  A long noncyclic path in
the refined tree is monotone by the reversal bound and projects, after
reinserting collapsed shift edges, to a long monotone path of $T_0$.
This proves (iv).

If there are no active colors, Lemma~\ref{gate:local-switch-tree} instead
gives a uniform bound on lengths of noncyclically stabilized paths.
The cyclic dictionary, Lemma~\ref{row:cyclic-dictionary}, then bounds
lengths of all nontrivially stabilized paths, so the refined action is
acylindrical.  The direct projection argument in
Proposition~\ref{transfer:ambient-qc}, applicable with the already established
hyperbolic rows and quasi-isometric edge maps, makes this a quasiconvex
hierarchy.  Its only potentially nonfree vertex is a hyperbolic root row,
which itself splits over a quasiconvex cyclic subgroup with free and
cyclic vertex groups.  The hierarchy theorem therefore makes $G$
virtually compact special in this case.
\end{proof}
